\chapter{モールス符号：\nt{GLUT}と\nt{GABA}のニューロンが話す言語について何がわかっているか}
\chaptermark{モールス符号}
\label{sec:code}

\section{記号が一つだけのアルファベット}

モールス符号には短点と長点（トンとツー）の二つの記号があり、その間に3種類の長さの休止がある。第\ref{sec:neurontypes}章で見たように、ニューロンのアルファベットはさらに貧しく、記号は一つしかない。スパイクは約1ミリ秒続き、1個のニューロンのスパイクは高さも形もすべて同じである。そのためメッセージのすべては休止に、つまり時刻の列$t_1,t_2,t_3,\dots$にある。長点のないモールス符号で、意味を運ぶのは短点のリズムだけである（図~\ref{fig:morse}）。

\begin{figure}[t]
\centering
\begin{tikzpicture}[>=Stealth,x=0.08cm,y=1cm]
% Morse letter: dot, dash, dot
\node[left,font=\small] at (-6,2.55) {モールス：};
\fill[black] (0,2.45) rectangle (6,2.65);
\fill[black] (14,2.45) rectangle (32,2.65);
\fill[black] (40,2.45) rectangle (46,2.65);
\node[right,font=\small,gray!40!black] at (52,2.55) {トン、ツー、トン：意味は長さと休止にある};
% stimulus onset
\node[left,font=\small] at (-6,0.4) {ニューロン：};
\draw[->,thick,green!45!black] (0,-0.55) -- (0,-0.05);
\node[below,font=\scriptsize,green!45!black] at (0,-0.55) {刺激};
\draw[gray!70!black] (0,0) -- (152,0) node[right,black] {\small $t$};
% spikes: single ones blue, the burst red
\foreach \s in {14,26,41,88,112,131}{\draw[very thick,blue!60!black] (\s,0) -- (\s,0.8);}
\foreach \s in {60,63,66,69}{\draw[very thick,red!70!black] (\s,0) -- (\s,0.8);}
% latency
\draw[<->,thick] (0,1.1) -- node[above,font=\scriptsize] {$t_1$：最初のスパイクの時刻} (14,1.1);
% burst
\draw[decorate,decoration={brace,amplitude=4pt},red!70!black] (58,0.95) -- (71,0.95) node[midway,above=4pt,font=\scriptsize] {バースト＝「ツー」};
% counting windows
\foreach \a/\b/\n in {0/50/3,50/100/5,100/150/2}{
  \draw[decorate,decoration={brace,amplitude=4pt,mirror},gray!50!black] (\a+1,-0.2) -- (\b-1,-0.2) node[midway,below=4pt,font=\small] {$n=\n$};}
\node[font=\scriptsize,gray!40!black] at (75,-1.05) {レート符号：各窓のスパイク数$n$};
\foreach \x in {0,50,100,150}{\draw[gray!60!black] (\x,-0.06) -- (\x,0.06);}
\end{tikzpicture}
\caption{モールス符号とニューロンの符号。同じスパイク列もいろいろに読める。$50\ms$の窓ごとにスパイクを数える（\ref{sec:rate}節）、遅延$t_1$~\eqref{eq:latency}を測る、あるいはバースト、つまり「ツー」~\eqref{eq:burst}に気づく。}
\label{fig:morse}
\end{figure}

そもそもどんな休止がありうるのか。下限は不応期で決まる。スパイクの後、ニューロンは約1ミリ秒の間は再び発火できないので、毎秒数百スパイクを超えることはない。上限は何もない。ニューロンは何分も黙っていられる。皮質の\nt{GLUT}ニューロンの発火率は毎秒1スパイク未満から数十スパイクの範囲にあり、ほとんどのニューロンはほとんどの時間を下限の近くで動作している。

第\ref{sec:glutcomputer}章と第\ref{sec:gaba}章では、スパイクで数を書き表すあれこれの方法を途中でその都度仮定してきた。ここで正面から問おう。\nt{GLUT}ニューロンと\nt{GABA}ニューロンは、互いにどんな言語で話しているのか。完全な答えはない。この言語は解読されていない。だが、わかっていることもいくつか確かにあり、そのほとんどは通信技術者のように考えれば導ける。通信路容量、雑音、受信機、伝送のコストである。

\section{どれだけ伝えられるか}

上限から始めよう。受け手がスパイクの時刻を精度$\Delta t$で区別するとする。$\Delta t$より小さなずれは受け手には区別できない。時間を不応期より短い長さ$\Delta t$のセルに刻み、各セルにはスパイクが1個あるかないかのどちらかだとする（図~\ref{fig:cells}）。平均発火率$r$では、スパイクがセルに入る確率は$p=r\Delta t$である。流れが最も多くの情報を運ぶのはセルが独立なときで、そのとき各セルはエントロピー$-p\log_2p-(1-p)\log_2(1-p)\approx p\log_2(e/p)$ビット（$p\ll1$）を与える。$\Delta t$で割ると、伝送速度の上限と、そのスパイク1個あたりの値が得られる~\cite{MacKay1952}。
\begin{empheq}[box=\keybox]{equation}
R_{\max} \;\approx\; r\,\log_2\frac{e}{r\,\Delta t}\ \ \text{bit/s},
\qquad
\frac{R_{\max}}{r} \;\approx\; \log_2\frac{e}{r\,\Delta t}\ \ \text{bit/スパイク}.
\label{eq:spikeentropy}
\end{empheq}
$r=10$スパイク/秒、$\Delta t=1\ms$では、スパイク1個あたり約8ビット、毎秒約80ビットである。

\begin{figure}[t]
\centering
\begin{tikzpicture}[>=Stealth,x=0.5cm,y=1cm]
% time cut into cells of width Delta t; a cell holds one impulse or none
\foreach \k in {0,...,23}{\draw[gray!55] (\k,0) rectangle (\k+1,0.7);}
\foreach \k in {2,7,8,15,21}{
  \fill[red!10] (\k,0) rectangle (\k+1,0.7);
  \draw[very thick,red!70!black] (\k+0.5,0.05) -- (\k+0.5,0.65);}
\foreach \k in {0,...,23}{
  \pgfmathtruncatemacro{\bit}{(\k==2)||(\k==7)||(\k==8)||(\k==15)||(\k==21)}
  \ifnum\bit=1 \def\bitcol{red!70!black}\else \def\bitcol{gray!60!black}\fi
  \node[font=\scriptsize,text=\bitcol] at (\k+0.5,-0.25) {\bit};}
\draw[<->,gray!60!black] (0,0.95) -- (1,0.95) node[midway,above,font=\scriptsize,black] {$\Delta t$};
\node[font=\small,anchor=east] at (-0.3,0.35) {スパイク};
\node[font=\small,anchor=east] at (-0.3,-0.25) {ビット};
\draw[decorate,decoration={brace,amplitude=5pt,mirror},gray!60!black] (0,-0.55) -- (24,-0.55)
  node[midway,below=6pt,font=\scriptsize,black,align=center] {数えるだけの受け手：「$n=5$」\\ セルが見える受け手：$24$個のうちどの$5$個か、つまり$\log_2\binom{24}{5}\approx15$ビット};
\end{tikzpicture}
\caption{容量~\eqref{eq:spikeentropy}はどこから来るか。時間を長さ$\Delta t$のセルに刻むと、各セルにはスパイクがあるか（$1$）ないか（$0$）のどちらかであり、スパイク列は2進の文字列になる。スパイクがセルに入る確率は$p=r\Delta t$である。セルが細かいほど文字列は長くなり、ビットも増える。スパイクを数えるだけのカウンタは、1が\emph{どこに}あるかに書かれたもののほとんどすべてを失う。}
\label{fig:cells}
\end{figure}

スパイク1個あたりのビット数は、時間精度に対数的にしか依存しない。精度$10\ms$でも、スパイク1個あたり約5ビットが残る。しかし受け手が1秒間のスパイクを数えるだけなら、$r=10$では1秒全体で2ビット未満しか得られない（\ref{sec:rate}節）。

\begin{keyblock}
\begin{proposition}[スパイク通信路の容量]
\label{prop:capacity}
平均発火率$r$のスパイク列を時間精度$\Delta t$で読むとき、運べる情報は~\eqref{eq:spikeentropy}を超えない。この容量のほとんどすべてはスパイクの\emph{時刻}にある。スパイクを数えるだけの受け手が同じ列から引き出すものは、時刻を1ミリ秒の精度で区別する受け手の数十分の1である。
\end{proposition}
\end{keyblock}

これは上限であって測定値ではない。刺激がわかっている感覚ニューロンでの測定では、スパイク1個あたり1ビットから数ビットであり、最もよい場合でその精度について計算した上限の半分ほどに達する~\cite{Rieke1997}。つまり、少なくとも神経系の入口では、スパイクの時刻は失われずに使われている。

\section{雑音のある通信路}

容量~\eqref{eq:spikeentropy}は雑音なしで計算した。雑音がどこで生じるかについては、三つの事実が確かにわかっている。

\emph{スパイク発生器そのものは精確である。}脳から組織片ごと取り出した皮質ニューロンに、時間的に速く変化する同じ電流を何度も与えると、スパイクは約1ミリ秒の精度で再現される~\cite{Mainen1995}。電流が一定なら、スパイクの時刻は試行ごとにばらけていく。精確な時刻を生むのはニューロン自体ではなく、入力の速い変化である。

\emph{シナプスは当てにならない。}\nt{GLUT}シナプスに達したスパイクは、ある確率$p$でしか神経伝達物質のひとまとまりを放出しない。海馬のシナプスでは半分以上のスパイクが受け手にまったく届かず、その原因は伝導の失敗ではなく、まさに放出のランダム性である~\cite{AllenStevens1994}。通信技術者にとって、これは消失通信路である。二つのニューロンの間にシナプスが複数あれば、ある程度は救われる。

\emph{生きた皮質では、スパイク列はランダムに見える。}スパイク間隔は、ランダムなポアソン過程とほぼ同じようにばらつく。同じ刺激を繰り返すと、窓内のスパイク数は分散が平均に近くなるように変動する~\cite{Softky1993}。

精確な素子がどうしてランダムな列を生むのか。ニューロンは数千の入力を受け取り、そのそれぞれが当てにならず、しかも第\ref{sec:gaba}章のように興奮と抑制が釣り合っている。膜電位は閾値付近で揺らぎ、スパイクの時刻はその揺らぎで決まる。これが雑音なのか、それとも私たちが読み方を知らないメッセージなのかは、大部分が未解決の問題である。「雑音」の一部はすでにメッセージだとわかった。視覚野では、ばらつきのかなりの部分が動物自身の運動に追従している~\cite{Stringer2019}。ここでの困難は原理的なものである。よく圧縮されたメッセージは、符号を知らない者にとってはランダムな列と区別がつかない。

\section{レート：遅いが確実}
\label{sec:rate}

最も単純な符号は\emph{レート符号}である。数は、ニューロンが窓$T$の間に出すスパイクの数に書かれている。消失も時刻の揺らぎも、この符号には怖くない。だが代償がある。ポアソン列なら、窓内のスパイク数$n$は平均$rT$、ばらつき$\sqrt{rT}$をもつ。
\begin{empheq}[box=\keybox]{equation}
\frac{\sigma_n}{\bar n} \;=\; \frac{1}{\sqrt{r\,T}} .
\label{eq:rateerror}
\end{empheq}
発火率を10パーセントの精度で知るには100個のスパイクが必要で、毎秒20スパイクなら5秒かかる。隣り合うレベルはばらつきの2倍ほど離れていれば区別できるので、発火率$r$までの範囲で時間$T$の間に区別できるレベルは約$\sqrt{rT}$個である。$r=10$、$T=1\,\text{s}$では3--4レベルで、2ビットに満たない。
脳はこんなに遅くは動かない。ヒトは一瞬だけ見せられた画像の中の動物を認識し、脳内の応答は約$150\ms$後に現れる~\cite{Thorpe1996}。粗く見積もって、この間に信号は約10段の逐次的な処理段を通り、1段あたり$10$--$20\ms$である。皮質のふつうの発火率では、各ニューロンは1段の間に0個か1個のスパイクしか出せず、その窓での発火率は定まらない。出口は二つある。多数のニューロンを使うか、時刻を見るかである。

\emph{多数のニューロン。}$N$個のニューロンが同じ数を運び、受け手がそれらのスパイクを足し合わせるとする。雑音が独立なら、式~\eqref{eq:rateerror}の$rT$は$NrT$に置き換わる。$r=20$の100個のニューロンは$15\ms$で30個のスパイクを出し、精度は約20パーセントになる。空間が時間の代わりをする。しかし隣り合う皮質ニューロンの雑音は独立ではない。隣どうしのペアのスパイク数の相関係数は、ふつう約$0.1$である~\cite{Zohary1994}。それが$c$なら、$N$個のニューロンの平均の分散は次のようになる。
\begin{empheq}[box=\keybox]{equation}
\sigma^2_N \;=\; \sigma^2_1\,\frac{1+(N-1)\,c}{N}
\;\xrightarrow[N\to\infty]{}\; c\,\sigma^2_1 .
\label{eq:correlated}
\end{empheq}

\begin{keyblock}
\begin{proposition}[平均化の限界]
\label{prop:pooling}
雑音が相関していると、群での平均化は誤差を$1/\sqrt{c}$分の1より小さくはできない。$c=0.1$ではこれは3分の1であり、この改善は数十個のニューロンですでに達成される。それ以上加えても無駄である。
\end{proposition}
\end{keyblock}

相関がすべて同じように有害なわけではない。精度を制限するのは、共通の雑音のうち\emph{信号そのものに似た}部分、つまり測定される量が変化したときと同じように群全体をずらす部分だけである~\cite{MorenoBote2014}。脳に実際にそのような雑音がどれだけあるかは、まだ調べられているところである。

多数のニューロンを使う方法はもう一つある。音の方向の検出器（図~\ref{fig:itd}）のように、\emph{どの}ニューロンが発火したかに数を書くことである。これは\emph{場所}による符号で、知られているうちで最も確実なものである。感覚ニューロンでは、配線の番号が皮膚のどの点に触れたか、音の高さはいくらかを告げ、これは何度も確かめられている。ニューロンの数の点では高価だが、速い。メッセージ1個に遅延1段で済む。

\section{時刻：速いが、どこから測るか}

二つ目の出口は、スパイクの時刻に数を書くことである。ニューロン~\eqref{eq:glutneuron}をとり、時刻$t=0$から、電圧をレベル$U$まで上げるような一定の入力を与える。電圧は$V=U\bigl(1-e^{-t/\tau}\bigr)$で増え、次の時刻に閾値$\theta$に達する。
\begin{empheq}[box=\keybox]{equation}
t_1 \;=\; \tau\,\ln\frac{U}{U-\theta}, \qquad U>\theta .
\label{eq:latency}
\end{empheq}
強い入力なら早いスパイク、弱い入力なら遅いスパイク、閾値以下ならスパイクなしである。$\tau=10\ms$では、入力$U=4\theta$で最初のスパイクは$2.9\ms$後、$U=2\theta$で$6.9\ms$後、$U=1.2\theta$で$17.9\ms$後に来る。たった1個のスパイクが数を運ぶ。

しかし$t_1$はどの時刻から測るのか。共通の時刻基準はなく、受け手は刺激がいつ始まったかを知らない。答えは三つある。

\emph{相対遅延。}二つのニューロンが同じ刺激を見るなら、その遅延の差$t_1^A-t_1^B$は入力だけで決まり、開始時刻にはよらない。時刻を比べるのは遅延つきの同時検出器である（図~\ref{fig:itd}）。$N$個のニューロンが1個ずつスパイクを出したなら、その到着の\emph{順序}そのものが最大$\log_2N!$ビットを運ぶ。10個のニューロンなら約22ビットであり、「発火したかどうか」という答えでは10ビット以下にしかならない。網膜では、出力ニューロンの最初のスパイクの相対遅延から画像を速く確実に復元できることが示されている~\cite{Gollisch2008}。

\emph{語頭としての一斉発火。}刺激の急な変化は、多数のニューロンをほぼ同時に発火させる。この一斉発火そのものが、モールス符号の語と語の間の長い休止のような開始の目印になり、受け手はそこから遅延を測る。

\emph{局所的なリズム。}第\ref{sec:gaba}章では、\nt{GLUT}--\nt{GABA}ループが周期$12$--$50\ms$の局所的なリズムを生んだ。これはクロックではないが、その1サイクルの中で、強い入力をもつニューロンは弱い入力をもつニューロンより早く発火でき、式~\eqref{eq:latency}は\emph{位相}による符号に変わる。数は、サイクルのどの部分でスパイクが来たかに書かれる。このような符号は観察されている。空間内の特定の場所を担当するニューロンは、動物が「自分の」場所を通過するにつれて、遅い局所リズムのサイクルのますます早い時点で発火する~\cite{OKeefe1993}。脳が位相をどれほど広く使っているかは、まだわかっていない。

\section{符号を選ぶのは受け手}

スパイク列には、発火率も、時刻も、順序もある。そのうちどれが読まれるかは受け手が決め、しかもただ一つのパラメータ、自分の時定数$\tau$で決める。

式~\eqref{eq:glutneuron}を時間で平均しよう。ニューロンに発火率$r_j$の独立な列が届くなら、平均電圧とその相対的な揺らぎは次のようになる。
\begin{empheq}[box=\keybox]{equation}
\bar V \;=\; \tau\sum_j w_j\,r_j ,
\qquad
\frac{\sigma_V}{\bar V} \;\approx\; \frac{1}{\sqrt{2\,r_{\text{in}}\,\tau}} ,
\label{eq:reader}
\end{empheq}
ここで2番目の式は重みが等しい場合のもので、$r_{\text{in}}$はすべての入力の発火率の合計である。$\tau$が大きいと、電圧はほとんど揺らがず、発火率の重みつき和に追従する。受け手は\emph{積分器}であり、発火率を読んで時刻を忘れる。$\tau$が小さいと、電圧はスパイクとスパイクの間に下がりきり、閾値に達するのは複数のスパイクが窓~\eqref{eq:window}の中に届いたときだけである。受け手は\emph{同時検出器}であり、時刻を読む（図~\ref{fig:reader}）。毎秒5スパイクの入力1000本、$\tau=10\ms$では揺らぎは約10パーセントで、ほぼ純粋な積分である。第\ref{sec:gaba}章のように抑制が$\tau$を$2\ms$に縮めると、揺らぎは20パーセント余りに増え、ニューロンは同時性に気づき始める。

\begin{figure}[t]
\centering
\begin{tikzpicture}[>=Stealth,x=0.115cm,y=1cm]
% common input: the same spike train for both receivers
\foreach \s in {6,21,22,38,52,55.5,59,62.5,66,69.5,73,76.5,80,83.5,87,90.5,94,97.5}{\draw[thick,gray!40!black] (\s,5.25) -- (\s,5.65);}
\draw[gray!70!black] (0,5.25) -- (105,5.25);
\node[left,font=\small] at (-1,5.45) {入力};
\node[above,font=\scriptsize,gray!40!black] at (13.5,5.65) {まばら};
\node[above,font=\scriptsize,gray!40!black] at (21.5,6.0) {ペア};
\node[above,font=\scriptsize,gray!40!black] at (75,5.65) {高頻度};
% axes and thresholds
\foreach \y/\th in {2.7/2.5,0/1.5}{
  \draw[gray!70!black] (0,\y) -- (106,\y) node[right,black] {\small $t$};
  \draw[densely dashed,gray!60!black] (0,\y+0.6*\th) -- (104,\y+0.6*\th) node[right,black,font=\small] {$\theta$};
}
\node[anchor=south west,font=\small] at (0,4.3) {$\tau=10\ms$：積分器};
\node[anchor=south east,font=\small] at (104,1.0) {$\tau=2\ms$：同時検出器};
% integrator: tau=10 ms, theta=2.5w
\begin{scope}[yshift=2.7cm]
\foreach \a/\b/\v in {0/6/0.000,6/21/1.000,21/22/1.223,22/38/2.107,38/52/1.425,52/55.5/1.351,55.5/59/1.952,59/62.5/2.376,62.5/66/0.000,66/69.5/1.000,69.5/73/1.705,73/76.5/2.201,76.5/80/0.000,80/83.5/1.000,83.5/87/1.705,87/90.5/2.201,90.5/94/0.000,94/97.5/1.000,97.5/104/1.705}{\draw[very thick,blue!60!black] plot[domain=\a:\b,samples=14] (\x,{0.6*\v*exp(-(\x-\a)/10)});}
\foreach \s/\p/\q in {6/0.000/1.000,21/0.223/1.223,22/1.107/2.107,38/0.425/1.425,52/0.351/1.351,55.5/0.952/1.952,59/1.376/2.376,62.5/1.674/2.500,62.5/2.500/0.000,66/0.000/1.000,69.5/0.705/1.705,73/1.201/2.201,76.5/1.551/2.500,76.5/2.500/0.000,80/0.000/1.000,83.5/0.705/1.705,87/1.201/2.201,90.5/1.551/2.500,90.5/2.500/0.000,94/0.000/1.000,97.5/0.705/1.705}{\draw[very thick,blue!60!black] (\s,0.6*\p) -- (\s,0.6*\q);}
\foreach \s in {62.5,76.5,90.5}{\draw[->,thick,red!70!black] (\s,1.58) -- (\s,1.95);}
\end{scope}
% coincidence: tau=2 ms, theta=1.5w
\begin{scope}[yshift=0cm]
\foreach \a/\b/\v in {0/6/0.000,6/21/1.000,21/22/1.001,22/38/0.000,38/52/1.000,52/55.5/1.001,55.5/59/1.174,59/62.5/1.204,62.5/66/1.209,66/69.5/1.210,69.5/73/1.210,73/76.5/1.210,76.5/80/1.210,80/83.5/1.210,83.5/87/1.210,87/90.5/1.210,90.5/94/1.210,94/97.5/1.210,97.5/104/1.210}{\draw[very thick,blue!60!black] plot[domain=\a:\b,samples=14] (\x,{0.6*\v*exp(-(\x-\a)/2)});}
\foreach \s/\p/\q in {6/0.000/1.000,21/0.001/1.001,22/0.607/1.500,22/1.500/0.000,38/0.000/1.000,52/0.001/1.001,55.5/0.174/1.174,59/0.204/1.204,62.5/0.209/1.209,66/0.210/1.210,69.5/0.210/1.210,73/0.210/1.210,76.5/0.210/1.210,80/0.210/1.210,83.5/0.210/1.210,87/0.210/1.210,90.5/0.210/1.210,94/0.210/1.210,97.5/0.210/1.210}{\draw[very thick,blue!60!black] (\s,0.6*\p) -- (\s,0.6*\q);}
\foreach \s in {22}{\draw[->,thick,red!70!black] (\s,0.98) -- (\s,1.35);}
\end{scope}
\foreach \x in {0,50,100}{\draw[gray!60!black] (\x,-0.06) -- (\x,0.06) node[below,font=\scriptsize] {\x\,ms};}
\end{tikzpicture}
\caption{同じ入力と二つの異なる受け手。各スパイクは電圧を$w$だけ上げ、その後はニューロンのモデル~\eqref{eq:glutneuron}のように電圧が流れ落ちる。赤い矢印は受け手のスパイクである。遅い受け手（$\tau=10\ms$、$\theta=2.5w$）はスパイクが\emph{多い}ところで発火し、ペアには気づかない。速い受け手（$\tau=2\ms$、$\theta=1.5w$）は窓~\eqref{eq:window}の中のペアにだけ発火し、高頻度でも同時でない列は見逃す。}
\label{fig:reader}
\end{figure}

測定によれば、活動している皮質では膜のコンダクタンスが静止時の数倍に増える~\cite{Destexhe2003}。つまりそこでは$\tau$が短く、動作中の皮質ニューロンは積分器よりも同時検出器に近い。

\begin{keyblock}
\begin{proposition}[符号を決めるのは受け手]
\label{prop:reader}
同じスパイク列が、遅い受け手にとっては発火率を、速い受け手にとっては時刻を、神経伝達物質が放出される体積にとっては数秒にわたって平滑化された濃度レベルを運ぶ（第\ref{sec:serocomputer}章）。どの符号が使われるかを決めるのは送り手ではなく受け手の時定数であり、それを抑制が調整する。
\end{proposition}
\end{keyblock}
\section{トンとツー：フィルタとしてのシナプス}

これまでシナプスは重みつきの配線にすぎなかった。実際には、シナプスの応答はそれ以前のスパイクに依存し、それがニューロンの言語に第二の記号を与える。

\emph{枯渇：シナプスは変化を伝える。}放出のたびに、放出準備のできた神経伝達物質の蓄えの割合$U$が使われ~\cite{Tsodyks1997}、蓄えは数百ミリ秒程度の時定数$\tau_{\text{rec}}$で回復する。発火率$r$では、スパイク1個への応答の振幅はほぼ$A(r)=A_0/(1+U r\tau_{\text{rec}})$に落ち着く。ここで$A_0$は休んだシナプスの応答である。受け手が単位時間に受け取る電流は次のようになる。
\begin{empheq}[box=\keybox]{equation}
r\,A(r) \;=\; \frac{A_0\,r}{1+U r\,\tau_{\text{rec}}}
\;\xrightarrow[r\to\infty]{}\; \frac{A_0}{U\tau_{\text{rec}}} .
\label{eq:depression}
\end{empheq}
$U=0.5$、$\tau_{\text{rec}}=0.8\,\text{s}$では、飽和は$1/(U\tau_{\text{rec}})=2.5$スパイク/秒からすでに始まる。発火率が5から20に増えても、定常電流は3分の1しか増えない。そのかわり、新しい発火率での最初のスパイクは蓄えがまだ以前のままのうちに届くので、電流は一瞬4倍に跳ね上がり、その後$\tau_{\text{rec}}/(1+Ur\tau_{\text{rec}})\approx90\ms$で減衰する（図~\ref{fig:depression}）。

このようなシナプスが伝えるのは発火率ではなくその\emph{変化}、本質的には相対変化$\Delta r/r$である~\cite{Abbott1997depression}。電子回路の技術者にとって、これは配線に直接組み込まれたハイパスフィルタである。

\begin{figure}[t]
\centering
\begin{tikzpicture}[>=Stealth,x=12.5cm,y=1cm]
% input rate
\draw[->,gray!70!black] (0,2.6) -- (0.9,2.6) node[right,black] {\small $t$};
\draw[->,gray!70!black] (0,2.6) -- (0,4.3) node[above,black] {\small $r$};
\draw[very thick,blue!60!black] (0,2.9) -- (0.2,2.9) -- (0.2,3.8) -- (0.8,3.8);
\node[left,font=\scriptsize] at (0,2.9) {$5$};
\node[left,font=\scriptsize] at (0,3.8) {$20$};
\node[right,font=\small,blue!60!black] at (0.4,4.1) {シナプスへの入力発火率};
% output current
\draw[->,gray!70!black] (0,0) -- (0.9,0) node[right,black] {\small $t$};
\draw[->,gray!70!black] (0,0) -- (0,2.45) node[left,black] {\small $rA$};
\draw[densely dashed,gray!60!black] (0,0.667) -- (0.8,0.667);
\draw[very thick,red!70!black] (0,0.5) -- (0.2,0.5) -- (0.2,2.0)
   plot[domain=0.2:0.8,samples=60] (\x,{0.667+(2.0-0.667)*exp(-(\x-0.2)/0.089)});
\node[left,font=\scriptsize] at (0,0.5) {$1.7$};
\node[left,font=\scriptsize] at (0,2.0) {$6.7$};
\node[right,font=\scriptsize] at (0.8,0.667) {$2.2$};
\node[right,font=\small,red!70!black] at (0.28,1.6) {受け手への電流：跳ね上がり、$\approx90\ms$で減衰};
\draw[gray!60!black] (0.2,-0.08) -- (0.2,0.08); \node[below,font=\scriptsize] at (0.2,0) {$0.2\,\text{s}$};
\draw[gray!60!black] (0.8,-0.08) -- (0.8,0.08); \node[below,font=\scriptsize] at (0.8,0) {$0.8\,\text{s}$};
\end{tikzpicture}
\caption{式~\eqref{eq:depression}による枯渇するシナプス、$U=0.5$、$\tau_{\text{rec}}=0.8\,\text{s}$。電流は毎秒$A_0$単位で表し、破線は定常電流である。定常電流は$1.7$から$2.2$に増えただけだが、跳躍の瞬間には$6.7$に達する。シナプスが受け手に伝えるのは、発火率がいくらかではなく、発火率が変わったということである。}
\label{fig:depression}
\end{figure}

\emph{促通：ツーとしてのバースト。}別のシナプスでは放出確率は小さいが、スパイクのたびに数十ミリ秒にわたって上がる。このようなシナプスでは、単発のスパイクはたいてい失われる。ところが\emph{バースト}、つまり数ミリ秒間隔の連続した数個のスパイクは、ほぼ確実に通過する。促通がなくても、バーストの$k$個のスパイクがすべて失われる確率は次のとおりである。
\begin{empheq}[box=\keybox]{equation}
P_{\text{loss}} \;=\; (1-p)^k ,
\label{eq:burst}
\end{empheq}
$p=0.2$では、単発のスパイクで$0.8$、4個のバーストで$0.41$である。促通はこの確率をさらに下げる。こうしてニューロンに第二の記号ができる。単発のスパイクがトン、バーストがツーである。枯渇するシナプスは休止の後の最初のスパイクを最もよく伝え、促通するシナプスはバーストを通して単発のトンを消す。バーストがより確実に伝わることは測定されている。脳が実際にそれを独立した記号として使っているというのは、もっともらしい仮説である~\cite{Lisman1997}。

\section{\nt{GABA}ニューロンはどう話すか}

皮質で最も数の多い\nt{GABA}ニューロンは速いタイプである~\cite{Tremblay2016}。そのスパイクは\nt{GLUT}ニューロンのものより短く、毎秒数百の頻度で一様に長時間出し続けることができ、ほとんど疲れない。そのシナプスはより確実である。各受け手との結合は多数の放出部位からなり、スパイクが失われることはほとんどない。その出力は受け手のスパイクが生まれる場所の近くにつくので、受け手が入力をどう足し合わせるかではなく、受け手が発火するかどうか、いつ発火するかを制御する。

速い\nt{GABA}ニューロン自身は近隣の多くの\nt{GLUT}ニューロンから興奮を受け取り、周囲の活動の総和を伝える。まさに第\ref{sec:gaba}章の割り算に必要なものである。さらに、速い\nt{GABA}ニューロンは互いに結合していて一緒に発火しやすい。前に述べた局所的なリズムを作っているのはこれらである~\cite{Cardin2009}。

モールス符号の言葉でいえば、\nt{GABA}ニューロンが受け持つのは文字ではなく\emph{休止}である。\nt{GLUT}スパイクの連続した流れを受け手が発火できる窓に区切り、同時性の窓を狭め、流れが大きくなりすぎると全体を抑える。

\begin{keyblock}
\begin{proposition}[役割の分担]
\label{prop:punctuation}
\nt{GLUT}ニューロンはメッセージの内容、つまり\emph{何を}と\emph{どれだけ}を運ぶ。速い\nt{GABA}ニューロンはその区切り、つまり\emph{いつ}話してよいか、\emph{どれだけの大きさで}話すかを運ぶ。もう一つのクラスは抑制するものを抑制することで、\emph{誰に対して}ゲートを開くかを決める。この分担の個々の部分は測定されている。一般規則としては作業仮説である。
\end{proposition}
\end{keyblock}

ほかに、出力が受け手の個々の入力につく、もっと遅い\nt{GABA}ニューロンもある。これは第\ref{sec:gaba}章の禁止~\eqref{eq:veto}として働く。ほかには触れずに、\nt{GLUT}スパイクの一つの流れを閉じる。

\nt{GABA}ニューロンを速いものと遅いものに分けるのは古い描像であり、不完全である。現在、皮質では少なくとも三つの大きなクラスが区別されており~\cite{Tremblay2016}、3番目のクラスは意外な構造をもつ。それは\nt{GLUT}ニューロンではなく別の\nt{GABA}ニューロン、とりわけ入力を閉じるものを抑制する~\cite{Pfeffer2013}。抑制の抑制は二重否定である。このようなニューロンは、受け手に興奮性のスパイクを一つも送らずにゲートを\emph{開く}。これをオンにするのは皮質の他の領域からの信号とブロードキャストチャネルなので、これはネットワークの一区画全体のイネーブル入力として働く。それが活動している間、禁止は解除され、流れが通る~\cite{Pi2013}。付録~\ref{sec:othermediators}では、この手法を脱抑制と呼んでいる。
\section{経済的な言語}

ニューロンの言語について確かにわかっている最後のことは、それが高価だということである。スパイクは、それが引き起こすシナプスの仕事と合わせて、約$10^9$個のATP分子を要し（げっ歯類の皮質での推定~\cite{Attwell2001}）、分子1個は約$10^{-19}$ジュールを与える。したがってスパイク1個のコストは$E_{\text{sp}}\sim10^{-10}$ジュールのオーダーである。脳全体は約$20$ワットを消費し、その半分、$P\sim10$ワットが信号の伝達に使われるとすると、ニューロン1個の平均発火率は次のようになる。
\begin{empheq}[box=\keybox]{equation}
\bar r \;\approx\; \frac{P}{N\,E_{\text{sp}}}
\;\sim\; \frac{10\ \text{W}}{8.6\cdot10^{10}\cdot10^{-10}\ \text{J}}
\;\approx\; 1\ \text{スパイク/秒}.
\label{eq:energy}
\end{empheq}
皮質についてのより詳しい推定では、さらに小さな値になる~\cite{Lennie2003}。脳の符号は\emph{スパース}でなければならない。速く動作できるのはニューロンのごく一部だけである。

式~\eqref{eq:spikeentropy}からわかるように、これはそれほど悪いことではない。精度$1\ms$では、毎秒100スパイクで動作するニューロンのスパイクは5ビット以下しか運ばないが、毎秒1スパイクで動作するニューロンのスパイクは最大11ビットを運ぶ。スパイクに代価を払うなら、まれに話すほうが得である。皮質は実際に精度をエネルギーと引き換えにする。食物が不足すると、ニューロンはスパイクの発火率を保ったままシナプス電流への支出を減らし、その応答は精度が下がる~\cite{Padamsey2022}。

モールス符号では、頻繁に使う文字は短い。英語の文章で最も頻繁な文字Eは短点1個である。ニューロンの符号は、わかっている限り、同じ規則に別の形で従っている。予想どおりのことにはスパイクをあまり使わない~\cite{Barlow1961}。一定の刺激のもとでニューロンは発火率を徐々に下げ、第\ref{sec:glutcomputer}章のセンサはレベルではなく変化に応答し、第\ref{sec:neurontypes}章の\nt{DOPA}ニューロンは報酬予測誤差だけを伝える。

\begin{keyblock}
\begin{proposition}[経済性]
\label{prop:economy}
エネルギーはニューロンの平均発火率を毎秒1スパイク程度に制限する。そのため\nt{GLUT}ニューロンの言語はスパースであり、スパイクをニュース、つまり変化したことや予測されなかったことに使う。エネルギーの限界は測定されている。脳の符号が1ジュールあたりのビット数を最大にするよう調整されているというのは、十分な根拠のある仮説である。
\end{proposition}
\end{keyblock}

\section{まとめ}

\begin{table}[t]
\centering
\footnotesize
\setlength{\tabcolsep}{4pt}
\renewcommand{\arraystretch}{1.15}
\begin{tabular}{>{\raggedright}p{2.6cm}>{\raggedright}p{2.3cm}>{\raggedright}p{3.0cm}>{\raggedright}p{2.0cm}>{\raggedright\arraybackslash}p{3.6cm}}
\toprule
書き方 & 運ぶもの & 読む者 & メッセージ1個の時間 & わかっていること \\
\midrule
発火したニューロンの番号（場所による符号） & どの値か & あらゆる受け手、同時検出器 & 遅延1段 & 感覚ニューロンと音の方向の検出で確立 \\
1個のニューロンの発火率 & 量 & $\tau$の大きい積分器、体積（第\ref{sec:serocomputer}章） & 数百ms--数秒 & 確立。$10$--$20\ms$の処理段には遅すぎる \\
群の発火率 & 量 & 数千の入力をもつニューロン & $10$--$50\ms$ & 確立。改善は相関~\eqref{eq:correlated}で制限される \\
最初のスパイクの時刻、順序 & 量、順序 & 遅延つきの同時検出器 & 遅延1段 & 神経系の入口で確立。皮質では議論中 \\
局所リズムの中の位相 & 量、位置 & 共通のリズムをもつニューロン & サイクル$12$--$50\ms$以上 & 一部の領域で観察。一般的な役割は仮説 \\
バースト（「ツー」） & 「重要」：確実な配送 & 促通するシナプス & $\sim10\ms$ & 確実性は測定済み。独立の記号というのは仮説 \\
発火率の変化 & ニュース & 枯渇するシナプス & $\sim0.1\,\text{s}$ & 確立 \\
速い\nt{GABA}ニューロンのスパイク & いつ、どれだけの大きさで & 近隣のすべての\nt{GLUT}ニューロン & $1$--$30\ms$ & リズムと割り算は測定済み。規則としての「句読点」は仮説 \\
\bottomrule
\end{tabular}
\caption{\nt{GLUT}ニューロンと\nt{GABA}ニューロンがスパイクでメッセージを書き表す方法。右の列は、測定されたことと推測されていることを区別している。}
\label{tab:code}
\end{table}

表~\ref{tab:code}は私たちが知っていることをまとめたものであり、知らないこともそこから見えてくる。皮質が群の発火率だけでなく、スパイクの精確な時刻をどれだけ使っているかはわからない。ニューロンに「単語」、つまり多数のニューロンのスパイクの安定した組み合わせで、全体として読まれるものがあるかどうかもわからない。そして、言語が脳の部位によって同じかどうかもわからない。モールス符号は一晩で覚えられる。その表がわかっているからである。ニューロンの言語の表はこれから作らなければならず、それを作るのはおそらく通信技術者だろう。

\section{問題}

\begin{exercise}[\lvlA]
\label{ex:04:1}
\emph{数値で見る容量。}$\Delta t=1\ms$とする。(a)~式~\eqref{eq:energy}のスパイクのコストのもとで、発火率毎秒$1$スパイクのニューロンは1ジュールあたり何ビットを与えるか。(b)~$r=10$で、1秒間のスパイクを数えるだけの受け手が引き出す情報は、上限~\eqref{eq:spikeentropy}の何分の1か。
\end{exercise}

\begin{exercise}[\lvlB]
\label{ex:04:2}
\emph{最適な発火率。}ニューロンは電力$P_0+rE_{\text{sp}}$を消費する。ここで$P_0$は$20$~Wの残り半分をすべてのニューロンで割ったもの、$E_{\text{sp}}=10^{-10}$~Jである。$\Delta t=1\ms$の式~\eqref{eq:spikeentropy}による1ジュールあたりのビット数$R_{\max}(r)/(P_0+rE_{\text{sp}})$が最大になる発火率$r$はいくらか。
\end{exercise}

\begin{exercise}[\lvlB]
\label{ex:04:3}
\emph{どの相関が有害か。}ニューロン$N=1000$個、分散$\sigma^2=1$、ペアの相関$c=0.1$。等しい傾き$f'=\mathbf 1$の場合と、$\mathbf 1^{\mathsf T}f'=0$となる$\pm1$の傾きの場合について、フィッシャー情報量$\mathcal I=f'^{\mathsf T}\Sigma^{-1}f'$（$\Sigma$は共分散行列）を計算せよ。両者は何倍異なるか。
\end{exercise}

\begin{exercise}[\lvlB]
\label{ex:04:4}
\emph{シミュレーション：同時検出器。}ニューロン~\eqref{eq:glutneuron}が毎秒$5$スパイクのポアソン入力を$1000$本受け取り、$w=1$である。閾値$\theta$は、自由な電圧がそれを毎秒1回横切るように選ぶ。$\tau=10$と$2\ms$のとき、ニューロンを確率$1/2$で発火させる同時入力の数$m$をシミュレーションで求めよ。
\end{exercise}

\begin{exercise}[\lvlB]
\label{ex:04:5}
\emph{設計：相対変化の検出器。}シナプス~\eqref{eq:depression}を設計せよ。毎秒$40$スパイクでの定常電流は毎秒$10$スパイクのときより$10\,\%$以上高くなってはならず、$40$への跳躍の後、電流は時定数$50\ms$以下で新しいレベルに落ち着かなければならない。最小の$\tau_{\text{rec}}$はいくらで、そのときの$U$はいくらか。
\end{exercise}

\begin{exercise}[\lvlA]
\label{ex:04:6}
\emph{最初のスパイクの時刻とバースト。}(a)~式~\eqref{eq:latency}から、最初のスパイクがちょうど時定数1個分の後に来る入力を求めよ。それは何に依存するか。(b)~放出確率$p=0.2$のとき、すべてが失われる確率を$5\,\%$未満にするには、バーストに何個のスパイクが必要か。
\end{exercise}

\begin{exercise}[\lvlC]
\label{ex:04:7}
\emph{研究：配置の中に何ビットあるか。}ニューロンを独立なセル~\eqref{eq:spikeentropy}とし、$r=10$、$\Delta t=1\ms$とする。(a)~$20$セルの「単語」のエントロピーを、スパイク数のエントロピーと残りに分解せよ。(b)~$N=30$と$10^4$個の記録から、単語の頻度によるエントロピーの推定をシミュレーションせよ。推定値はどれだけ過小になるか。
\end{exercise}
